%% kapa: exam=1 round=20 year=2009 period=2 form=A subject=appraisal_law subject_ko=부동산관계법규 seq=38 number=78
%% source: Q-Net 감정평가사 기출문제 (cst00309) artlSeq=5024178 file=extracted/1008/5024178/2009년 제20회 감정평가사 1차시험 기출문제/2교시.hwp
%% conversion: kordoc 4.24.51 -> tools/gen1.py (auto); review=proofread-sonnet; flags=none
\begin{question}{78}
부동산등기법령 및 판례에 의할 때 등기의 유효요건에 관한 설명으로 옳은 것은?
\begin{choices}
  \item 등기가 불법 말소된 후 제3자 명의의 등기가 행하여진 경우에는 말소회복등기를 할 수 없다.
  \item 동일 부동산에 관하여 동일한 명의로 중복하여 보존등기가 행하여진 경우에는 먼저 이루어진 등기가 원인무효가 아닌 한 나중에 이루어진 등기를 무효로 한다.
  \item 중간생략등기는 부동산등기특별조치법에 반하는 것으로 사법상의 효력이 인정되지 않는다.
  \item 소유권이전등기 원인이 증여인데도 이를 매매로 기재한 경우에는 그 등기는 무효가 된다.
  \item 등기부의 전부 또는 일부가 멸실한 경우에는 등기부상의 권리가 모두 소멸한다.
\end{choices}
\end{question}
\answer{2}
